% ============================================================
% Figura: Processo metodológico
% Requer no preâmbulo:
%   \usepackage{tikz}
%   \usetikzlibrary{positioning,arrows.meta,fit,backgrounds,calc}
% Uso no corpo do texto:  % ============================================================
% Figura: Processo metodológico
% Requer no preâmbulo:
%   \usepackage{tikz}
%   \usetikzlibrary{positioning,arrows.meta,fit,backgrounds,calc}
% Uso no corpo do texto:  % ============================================================
% Figura: Processo metodológico
% Requer no preâmbulo:
%   \usepackage{tikz}
%   \usetikzlibrary{positioning,arrows.meta,fit,backgrounds,calc}
% Uso no corpo do texto:  % ============================================================
% Figura: Processo metodológico
% Requer no preâmbulo:
%   \usepackage{tikz}
%   \usetikzlibrary{positioning,arrows.meta,fit,backgrounds,calc}
% Uso no corpo do texto:  \input{figura-metodologia}
% ============================================================
\begin{figure}[t]
\centering
\resizebox{\textwidth}{!}{%
\begin{tikzpicture}[
  node distance=5mm and 12mm,
  etapa/.style={rectangle, rounded corners=2pt, draw=black!70, fill=white,
                text width=5.2cm, align=center, font=\small,
                minimum height=1.2cm, inner sep=4pt},
  apoio/.style={etapa, dashed, fill=black!5},
  seta/.style={-{Stealth[length=2.5mm]}, thick, draw=black!70},
  fase/.style={draw=black!40, fill=black!4, rounded corners=4pt, inner sep=3mm},
  titulo/.style={font=\bfseries\small, minimum height=9mm, anchor=north, text width=5.2cm, align=center}
]

% ---------------- Fase 1: Levantamento ----------------
\node[titulo] (t1) {Fase 1\\Levantamento};
\node[etapa, below=of t1] (e1)
  {Identificação dos portais de prefeituras e câmaras\\{\footnotesize(Radar da Transparência; 314 URLs)}};
\node[etapa, below=of e1] (e2)
  {Localização manual dos canais de ouvidoria\\{\footnotesize(142 municípios)}};
\node[etapa, below=of e2] (e3)
  {Classificação do tipo de ouvidoria\\[2pt]
   {\scriptsize integrada ao portal $\cdot$ sistema próprio (domínio/subdomínio)
    $\cdot$ terceirizada $\cdot$ Fala.BR
    $\cdot$ não localizada $\cdot$ indisponível}};
    
% ---------------- Fase 2: Avaliação ----------------
\node[titulo, right=of t1] (t2) {Fase 2\\Avaliação automatizada};
\node[apoio, below=of t2] (e4)
  {\textit{Scripts} em Python\\{\footnotesize(requisições diretas ao serviço de avaliação)}};
\node[etapa, below=of e4] (e5)
  {Avaliação com AMAWeb e AccessMonitor\\{\footnotesize (WCAG 2.1) prefeituras, câmaras e ouvidorias; até cinco tentativas}};
\node[apoio, below=of e5] (e5b)
  {Contingência: envio do HTML baixado\\{\footnotesize(páginas bloqueadas por serviços de proteção)}};
\node[etapa, below=of e5b] (e6)
  {Registro de evidências e validação\\{\footnotesize JSON bruto, PDF e \textit{hash}; uma avaliação por município, ambiente e ferramenta}};

% ---------------- Fase 3: Consolidação e análise ----------------
\node[titulo, right=of t2] (t3) {Fase 3\\Consolidação e análise};
\node[etapa, below=of t3] (e7)
  {Planilha consolidada\\{\footnotesize URL, ferramenta, situação, nota, falhas, data/hora e \textit{hash}}};
\node[apoio, below=of e7] (e8)
  {Dados complementares\\{\footnotesize População (IBGE, 2026), PIB \textit{per capita} (2023), IDHM (2010)}};
\node[etapa, below=of e8] (e9)
  {Estatística descritiva (QP1)\\{\footnotesize média, mediana, dispersão e faixas de nota}};
\node[etapa, below=of e9] (e10)
  {Comparação entre ambientes e por tipo de ouvidoria (QP2 e QP3)\\{\footnotesize diferenças prefeitura $-$ ouvidoria e prefeitura $-$ câmara; Wilcoxon pareado}};
\node[etapa, below=of e10] (e11)
  {Correlação com indicadores (QP4), erros e ranqueamento\\{\footnotesize Spearman; erros mais frequentes; média das três notas}};

% ---------------- Caixas das fases (mesma altura) ----------------
\begin{scope}[on background layer]
  \node[fase, fit=(t1)(e1)(e3)(e3|-e11.south)] (f1) {};
  \node[fase, fit=(t2)(e4)(e6)(e6|-e11.south)] (f2) {};
  \node[fase, fit=(t3)(e7)(e11)] (f3) {};
\end{scope}

% ---------------- Setas ----------------
\draw[seta] (e1) -- (e2);
\draw[seta] (e2) -- (e3);
\draw[seta] (e4) -- (e5);
\draw[seta] (e5) -- (e5b);
\draw[seta] (e5b) -- (e6);
\draw[seta] (e7) -- (e8);
\draw[seta] (e8) -- (e9);
\draw[seta] (e9) -- (e10);
\draw[seta] (e10) -- (e11);

\draw[seta] (f1.east) -- node[above, font=\scriptsize, align=center]{URLs}   (f2.west);
\draw[seta] (f2.east) -- node[above, font=\scriptsize, align=center]{notas}  (f3.west);

\end{tikzpicture}%
}
\caption{Processo metodológico adotado na avaliação da acessibilidade dos portais das prefeituras, dos portais das câmaras e dos canais de ouvidoria dos municípios de Mato Grosso.}
\label{fig:metodologia}
\end{figure}
% ============================================================
\begin{figure}[t]
\centering
\resizebox{\textwidth}{!}{%
\begin{tikzpicture}[
  node distance=5mm and 12mm,
  etapa/.style={rectangle, rounded corners=2pt, draw=black!70, fill=white,
                text width=5.2cm, align=center, font=\small,
                minimum height=1.2cm, inner sep=4pt},
  apoio/.style={etapa, dashed, fill=black!5},
  seta/.style={-{Stealth[length=2.5mm]}, thick, draw=black!70},
  fase/.style={draw=black!40, fill=black!4, rounded corners=4pt, inner sep=3mm},
  titulo/.style={font=\bfseries\small, minimum height=9mm, anchor=north, text width=5.2cm, align=center}
]

% ---------------- Fase 1: Levantamento ----------------
\node[titulo] (t1) {Fase 1\\Levantamento};
\node[etapa, below=of t1] (e1)
  {Identificação dos portais de prefeituras e câmaras\\{\footnotesize(Radar da Transparência; 314 URLs)}};
\node[etapa, below=of e1] (e2)
  {Localização manual dos canais de ouvidoria\\{\footnotesize(142 municípios)}};
\node[etapa, below=of e2] (e3)
  {Classificação do tipo de ouvidoria\\[2pt]
   {\scriptsize integrada ao portal $\cdot$ sistema próprio (domínio/subdomínio)
    $\cdot$ terceirizada $\cdot$ Fala.BR
    $\cdot$ não localizada $\cdot$ indisponível}};
    
% ---------------- Fase 2: Avaliação ----------------
\node[titulo, right=of t1] (t2) {Fase 2\\Avaliação automatizada};
\node[apoio, below=of t2] (e4)
  {\textit{Scripts} em Python\\{\footnotesize(requisições diretas ao serviço de avaliação)}};
\node[etapa, below=of e4] (e5)
  {Avaliação com AMAWeb e AccessMonitor\\{\footnotesize (WCAG 2.1) prefeituras, câmaras e ouvidorias; até cinco tentativas}};
\node[apoio, below=of e5] (e5b)
  {Contingência: envio do HTML baixado\\{\footnotesize(páginas bloqueadas por serviços de proteção)}};
\node[etapa, below=of e5b] (e6)
  {Registro de evidências e validação\\{\footnotesize JSON bruto, PDF e \textit{hash}; uma avaliação por município, ambiente e ferramenta}};

% ---------------- Fase 3: Consolidação e análise ----------------
\node[titulo, right=of t2] (t3) {Fase 3\\Consolidação e análise};
\node[etapa, below=of t3] (e7)
  {Planilha consolidada\\{\footnotesize URL, ferramenta, situação, nota, falhas, data/hora e \textit{hash}}};
\node[apoio, below=of e7] (e8)
  {Dados complementares\\{\footnotesize População (IBGE, 2026), PIB \textit{per capita} (2023), IDHM (2010)}};
\node[etapa, below=of e8] (e9)
  {Estatística descritiva (QP1)\\{\footnotesize média, mediana, dispersão e faixas de nota}};
\node[etapa, below=of e9] (e10)
  {Comparação entre ambientes e por tipo de ouvidoria (QP2 e QP3)\\{\footnotesize diferenças prefeitura $-$ ouvidoria e prefeitura $-$ câmara; Wilcoxon pareado}};
\node[etapa, below=of e10] (e11)
  {Correlação com indicadores (QP4), erros e ranqueamento\\{\footnotesize Spearman; erros mais frequentes; média das três notas}};

% ---------------- Caixas das fases (mesma altura) ----------------
\begin{scope}[on background layer]
  \node[fase, fit=(t1)(e1)(e3)(e3|-e11.south)] (f1) {};
  \node[fase, fit=(t2)(e4)(e6)(e6|-e11.south)] (f2) {};
  \node[fase, fit=(t3)(e7)(e11)] (f3) {};
\end{scope}

% ---------------- Setas ----------------
\draw[seta] (e1) -- (e2);
\draw[seta] (e2) -- (e3);
\draw[seta] (e4) -- (e5);
\draw[seta] (e5) -- (e5b);
\draw[seta] (e5b) -- (e6);
\draw[seta] (e7) -- (e8);
\draw[seta] (e8) -- (e9);
\draw[seta] (e9) -- (e10);
\draw[seta] (e10) -- (e11);

\draw[seta] (f1.east) -- node[above, font=\scriptsize, align=center]{URLs}   (f2.west);
\draw[seta] (f2.east) -- node[above, font=\scriptsize, align=center]{notas}  (f3.west);

\end{tikzpicture}%
}
\caption{Processo metodológico adotado na avaliação da acessibilidade dos portais das prefeituras, dos portais das câmaras e dos canais de ouvidoria dos municípios de Mato Grosso.}
\label{fig:metodologia}
\end{figure}
% ============================================================
\begin{figure}[t]
\centering
\resizebox{\textwidth}{!}{%
\begin{tikzpicture}[
  node distance=5mm and 12mm,
  etapa/.style={rectangle, rounded corners=2pt, draw=black!70, fill=white,
                text width=5.2cm, align=center, font=\small,
                minimum height=1.2cm, inner sep=4pt},
  apoio/.style={etapa, dashed, fill=black!5},
  seta/.style={-{Stealth[length=2.5mm]}, thick, draw=black!70},
  fase/.style={draw=black!40, fill=black!4, rounded corners=4pt, inner sep=3mm},
  titulo/.style={font=\bfseries\small, minimum height=9mm, anchor=north, text width=5.2cm, align=center}
]

% ---------------- Fase 1: Levantamento ----------------
\node[titulo] (t1) {Fase 1\\Levantamento};
\node[etapa, below=of t1] (e1)
  {Identificação dos portais de prefeituras e câmaras\\{\footnotesize(Radar da Transparência; 314 URLs)}};
\node[etapa, below=of e1] (e2)
  {Localização manual dos canais de ouvidoria\\{\footnotesize(142 municípios)}};
\node[etapa, below=of e2] (e3)
  {Classificação do tipo de ouvidoria\\[2pt]
   {\scriptsize integrada ao portal $\cdot$ sistema próprio (domínio/subdomínio)
    $\cdot$ terceirizada $\cdot$ Fala.BR
    $\cdot$ não localizada $\cdot$ indisponível}};
    
% ---------------- Fase 2: Avaliação ----------------
\node[titulo, right=of t1] (t2) {Fase 2\\Avaliação automatizada};
\node[apoio, below=of t2] (e4)
  {\textit{Scripts} em Python\\{\footnotesize(requisições diretas ao serviço de avaliação)}};
\node[etapa, below=of e4] (e5)
  {Avaliação com AMAWeb e AccessMonitor\\{\footnotesize (WCAG 2.1) prefeituras, câmaras e ouvidorias; até cinco tentativas}};
\node[apoio, below=of e5] (e5b)
  {Contingência: envio do HTML baixado\\{\footnotesize(páginas bloqueadas por serviços de proteção)}};
\node[etapa, below=of e5b] (e6)
  {Registro de evidências e validação\\{\footnotesize JSON bruto, PDF e \textit{hash}; uma avaliação por município, ambiente e ferramenta}};

% ---------------- Fase 3: Consolidação e análise ----------------
\node[titulo, right=of t2] (t3) {Fase 3\\Consolidação e análise};
\node[etapa, below=of t3] (e7)
  {Planilha consolidada\\{\footnotesize URL, ferramenta, situação, nota, falhas, data/hora e \textit{hash}}};
\node[apoio, below=of e7] (e8)
  {Dados complementares\\{\footnotesize População (IBGE, 2026), PIB \textit{per capita} (2023), IDHM (2010)}};
\node[etapa, below=of e8] (e9)
  {Estatística descritiva (QP1)\\{\footnotesize média, mediana, dispersão e faixas de nota}};
\node[etapa, below=of e9] (e10)
  {Comparação entre ambientes e por tipo de ouvidoria (QP2 e QP3)\\{\footnotesize diferenças prefeitura $-$ ouvidoria e prefeitura $-$ câmara; Wilcoxon pareado}};
\node[etapa, below=of e10] (e11)
  {Correlação com indicadores (QP4), erros e ranqueamento\\{\footnotesize Spearman; erros mais frequentes; média das três notas}};

% ---------------- Caixas das fases (mesma altura) ----------------
\begin{scope}[on background layer]
  \node[fase, fit=(t1)(e1)(e3)(e3|-e11.south)] (f1) {};
  \node[fase, fit=(t2)(e4)(e6)(e6|-e11.south)] (f2) {};
  \node[fase, fit=(t3)(e7)(e11)] (f3) {};
\end{scope}

% ---------------- Setas ----------------
\draw[seta] (e1) -- (e2);
\draw[seta] (e2) -- (e3);
\draw[seta] (e4) -- (e5);
\draw[seta] (e5) -- (e5b);
\draw[seta] (e5b) -- (e6);
\draw[seta] (e7) -- (e8);
\draw[seta] (e8) -- (e9);
\draw[seta] (e9) -- (e10);
\draw[seta] (e10) -- (e11);

\draw[seta] (f1.east) -- node[above, font=\scriptsize, align=center]{URLs}   (f2.west);
\draw[seta] (f2.east) -- node[above, font=\scriptsize, align=center]{notas}  (f3.west);

\end{tikzpicture}%
}
\caption{Processo metodológico adotado na avaliação da acessibilidade dos portais das prefeituras, dos portais das câmaras e dos canais de ouvidoria dos municípios de Mato Grosso.}
\label{fig:metodologia}
\end{figure}
% ============================================================
\begin{figure}[t]
\centering
\resizebox{\textwidth}{!}{%
\begin{tikzpicture}[
  node distance=5mm and 12mm,
  etapa/.style={rectangle, rounded corners=2pt, draw=black!70, fill=white,
                text width=5.2cm, align=center, font=\small,
                minimum height=1.2cm, inner sep=4pt},
  apoio/.style={etapa, dashed, fill=black!5},
  seta/.style={-{Stealth[length=2.5mm]}, thick, draw=black!70},
  fase/.style={draw=black!40, fill=black!4, rounded corners=4pt, inner sep=3mm},
  titulo/.style={font=\bfseries\small, minimum height=9mm, anchor=north, text width=5.2cm, align=center}
]

% ---------------- Fase 1: Levantamento ----------------
\node[titulo] (t1) {Fase 1\\Levantamento};
\node[etapa, below=of t1] (e1)
  {Identificação dos portais de prefeituras e câmaras\\{\footnotesize(Radar da Transparência; 314 URLs)}};
\node[etapa, below=of e1] (e2)
  {Localização manual dos canais de ouvidoria\\{\footnotesize(142 municípios)}};
\node[etapa, below=of e2] (e3)
  {Classificação do tipo de ouvidoria\\[2pt]
   {\scriptsize integrada ao portal $\cdot$ sistema próprio (domínio/subdomínio)
    $\cdot$ terceirizada $\cdot$ Fala.BR
    $\cdot$ não localizada $\cdot$ indisponível}};
    
% ---------------- Fase 2: Avaliação ----------------
\node[titulo, right=of t1] (t2) {Fase 2\\Avaliação automatizada};
\node[apoio, below=of t2] (e4)
  {\textit{Scripts} em Python\\{\footnotesize(requisições diretas ao serviço de avaliação)}};
\node[etapa, below=of e4] (e5)
  {Avaliação com AMAWeb e AccessMonitor\\{\footnotesize (WCAG 2.1) prefeituras, câmaras e ouvidorias; até cinco tentativas}};
\node[apoio, below=of e5] (e5b)
  {Contingência: envio do HTML baixado\\{\footnotesize(páginas bloqueadas por serviços de proteção)}};
\node[etapa, below=of e5b] (e6)
  {Registro de evidências e validação\\{\footnotesize JSON bruto, PDF e \textit{hash}; uma avaliação por município, ambiente e ferramenta}};

% ---------------- Fase 3: Consolidação e análise ----------------
\node[titulo, right=of t2] (t3) {Fase 3\\Consolidação e análise};
\node[etapa, below=of t3] (e7)
  {Planilha consolidada\\{\footnotesize URL, ferramenta, situação, nota, falhas, data/hora e \textit{hash}}};
\node[apoio, below=of e7] (e8)
  {Dados complementares\\{\footnotesize População (IBGE, 2026), PIB \textit{per capita} (2023), IDHM (2010)}};
\node[etapa, below=of e8] (e9)
  {Estatística descritiva (QP1)\\{\footnotesize média, mediana, dispersão e faixas de nota}};
\node[etapa, below=of e9] (e10)
  {Comparação entre ambientes e por tipo de ouvidoria (QP2 e QP3)\\{\footnotesize diferenças prefeitura $-$ ouvidoria e prefeitura $-$ câmara; Wilcoxon pareado}};
\node[etapa, below=of e10] (e11)
  {Correlação com indicadores (QP4), erros e ranqueamento\\{\footnotesize Spearman; erros mais frequentes; média das três notas}};

% ---------------- Caixas das fases (mesma altura) ----------------
\begin{scope}[on background layer]
  \node[fase, fit=(t1)(e1)(e3)(e3|-e11.south)] (f1) {};
  \node[fase, fit=(t2)(e4)(e6)(e6|-e11.south)] (f2) {};
  \node[fase, fit=(t3)(e7)(e11)] (f3) {};
\end{scope}

% ---------------- Setas ----------------
\draw[seta] (e1) -- (e2);
\draw[seta] (e2) -- (e3);
\draw[seta] (e4) -- (e5);
\draw[seta] (e5) -- (e5b);
\draw[seta] (e5b) -- (e6);
\draw[seta] (e7) -- (e8);
\draw[seta] (e8) -- (e9);
\draw[seta] (e9) -- (e10);
\draw[seta] (e10) -- (e11);

\draw[seta] (f1.east) -- node[above, font=\scriptsize, align=center]{URLs}   (f2.west);
\draw[seta] (f2.east) -- node[above, font=\scriptsize, align=center]{notas}  (f3.west);

\end{tikzpicture}%
}
\caption{Processo metodológico adotado na avaliação da acessibilidade dos portais das prefeituras, dos portais das câmaras e dos canais de ouvidoria dos municípios de Mato Grosso.}
\label{fig:metodologia}
\end{figure}